% !TEX root = ../../main.tex
\subsection{方案安全性分析}
\label{subsec:security_analysis}

本节对 FoodShield 的安全设计进行半形式化分析，依次针对 3.4 节所列的八项安全目标，说明系统采用的对应机制如何在威胁模型假设下提供安全保障。

\subsubsection*{目标1：匿名性}

\textbf{目标}：骑手端和外部攻击者在日常通信中无法从系统返回的数据中确定用户真实身份或将多次通信关联至同一用户。

\textbf{分析}：系统为每位用户生成 PID$= \operatorname{HMAC\text{-}SM3}(K_{\mathrm{master}},\allowbreak \mathit{userID} \| r)$。在$K_{\mathrm{master}}$保密的前提下，HMAC-SM3 具有伪随机性，使得 PID 在计算上不可区分于随机字符串，攻击者无法从 PID 还原$\mathit{userID}$。骑手端接口响应不暴露 PID 字段，管理员默认视图不展示 PID 与用户名的映射，普通业务流程中没有任何接口将$(\mathrm{PID}, \mathit{username})$同时返回给骑手或外部调用方。

\textbf{复杂度估计}：在不知道$K_{\mathrm{master}}$的情况下，穷举攻击 HMAC-SM3 需要$2^{256}$次查询，超出实际计算能力。

\subsubsection*{目标2：认证性}

\textbf{目标}：只有持有合法 Token 的订单参与方才能进入对应订单通信通道，攻击者无法伪造或复用 Token 冒充合法用户。

\textbf{分析}：Token 公式为$\operatorname{HMAC\text{-}SM3}(K_{\mathrm{order}},\allowbreak \mathit{orderID} \| \mathrm{PID} \| \mathit{timestamp} \| \mathit{nonce})$。在$K_{\mathrm{order}}$保密的前提下，HMAC 的不可伪造性保证攻击者无法在不知道密钥的情况下构造任何有效 Token。时间戳有效期（默认 3600 秒）限制了截获 Token 后的可利用时间窗口；nonce 的引入确保即使时间戳相同，重放攻击也无法复现合法 Token。

\textbf{复杂度估计}：在随机预言模型下，HMAC-SM3 满足 EUF-CMA（存在不可伪造性），伪造成功概率不超过$q/2^{256}$，其中$q$为查询次数。

\subsubsection*{目标3：机密性}

\textbf{目标}：消息内容在数据库泄露场景下不被直接读取；管理员无需接触明文即可完成审计操作。

\textbf{分析}：消息明文经 SM4-CBC 加密后以密文形式存储，每条消息使用独立随机 IV，相同明文在不同时间产生不同密文，防止基于密文相等进行的内容推断。管理员消息查询接口只返回\texttt{message\_hash}和元数据字段，不返回\texttt{content}字段，在接口层面实现了审计与明文访问的权限分离。

\textbf{SM4-CBC 安全性}：SM4 是国家密码局发布的分组密码标准，密钥长度 128 位，CBC 模式在 IV 随机且保密的条件下满足 IND-CPA 安全性，穷举密钥需要$2^{128}$次操作。

\subsubsection*{目标4：完整性}

\textbf{目标}：任何对历史通信记录（消息内容、顺序、数量）的修改均可被检测。

\textbf{分析}：每条消息的 \texttt{message\_hash} 绑定了订单号、PID、角色、明文内容和时间戳五个字段，SM3 的抗碰撞性保证修改任意字段后哈希值以压倒性概率改变（碰撞概率不超过$2^{-128}$）。Merkle Tree 将所有消息哈希组织为树形结构，任意叶子节点变化均会级联改变 Root 值。三端摘要一致性机制进一步将 Root 快照分散至用户端、骑手端和服务端三方留存，攻击者需同时篡改三方数据才能伪造一致结果。

\subsubsection*{目标5：可追责性与可审计性}

\textbf{目标}：发生纠纷时能够通过合法流程追溯违规行为；管理员操作本身也处于可追踪的审计视野中。

\textbf{分析}：条件溯源的“完整性前置 + 关键词双重门槛”设计使溯源操作与违规证据直接绑定，防止无故溯源；所有管理员操作（包括溯源成功、溯源失败、完整性验证结果）均写入 \texttt{audit\_logs}，形成不可篡改的操作时间线（在 Merkle 保护范围内）。

\subsubsection*{目标6：抗滥用能力}

\textbf{目标}：防止机器批量注册生成大量匿名身份，以及利用匿名性发起垃圾消息或骚扰。

\textbf{分析}：注册阶段的数学验证码（服务端 session 一次性）要求攻击者每次注册都实时解析新题目，不能通过重放答案完成批量注册；口令长度下限和用户名唯一性约束进一步增加了批量注册的操作成本；订单号与 PID 的绑定机制使每次通信都要经过 Token 认证，无状态连接无法进入通信通道。

\subsubsection*{目标7：隐私最小化}

\textbf{目标}：系统各端只获取完成自身功能所必需的最小信息集合。

\textbf{分析}：骑手端接口不返回用户 PID 和备注字段；管理员默认接口不返回消息密文和用户真实身份；\texttt{orders.note} 备注字段在骑手端和管理员端的 API 响应中被屏蔽；\texttt{users.pid\_r}（随机数）不通过任何前端接口暴露。字段最小化从数据结构和接口设计两个层面共同实现隐私最小化原则。

\subsubsection*{综合评估}

综合来看，FoodShield 通过 HMAC-SM3 订单认证、SM4-CBC 加密存储、SM3 消息摘要、Merkle Tree 完整性验证、条件溯源约束和三端摘要一致性增强等机制的组合，在外卖订单通信场景中实现了匿名性、认证性、机密性、完整性、可追责性、可审计性和抗滥用能力的多目标安全保障。系统的安全设计以订单为核心索引，将业务流程与密码学机制紧密耦合，避免了安全机制游离于业务逻辑之外导致的防护盲区。
